\chapter{Framework Evolution and Future Developments}
\label{chap:future-developments}

The framework developed in this thesis was intentionally designed as a
lightweight layer on top of Kathar\'a. Its current scope is therefore limited
and explicit: it provides a declarative representation of a laboratory,
validates that representation, generates standard Kathar\'a artefacts, assists
with laboratory lifecycle operations, and supports a set of repeatable
measurements. The experimental evaluation in
\Cref{chap:experimental-evaluation} shows that this foundation is already
sufficient to build and observe a non-trivial routed environment under
substantially different operating conditions. At the same time, the same
experiment also identifies several directions in which the framework could be
generalised.

The purpose of this chapter is not to describe functionality that is already
implemented. Instead, it presents a possible evolution path from the current
prototype towards a broader framework for network experimentation. The main
idea is to preserve the characteristics that motivated the original design---a
readable declarative configuration, explicit validation, standard Kathar\'a
output, and simple command-line operation---while extending the framework with
reusable experiment definitions, runtime verification, richer experiment
management, and a cleaner separation between the logical network model and the
execution backend. These developments are discussed as future work and should
therefore be interpreted as architectural proposals rather than as features of
the present implementation.

\section{Extension to Network Attack Testing}
\label{sec:future-attack-testing}

The controlled scenario presented in the previous chapter demonstrates that the
framework can support more than topology generation. A generated laboratory can
also host application services, execute coordinated client measurements, apply
a bounded adverse workload, and collect observations that distinguish network
reachability from application-level availability. In the current repository,
however, this workflow is tied to a specific experiment and to a set of scripts
written for that experiment. A natural evolution would be to extract the
reusable concepts from that implementation and introduce a generic subsystem
for controlled security and resilience testing.

The objective of such an extension would not be to turn the project into a
general-purpose attack tool. Instead, it would support bounded experiments that
study denial-of-service and related resilience conditions inside an isolated
environment \cite{rfc4732}. The value of Kathar\'a in this context is precisely
that potentially disruptive network conditions can be reproduced inside an
isolated and explicitly configured environment. The future framework should
therefore model adverse conditions as \emph{experimental scenarios}: bounded,
repeatable changes applied to nodes that belong to the current laboratory and
observed through a predefined measurement procedure. This interpretation also
broadens the scope beyond malicious traffic. From the point of view of the
framework, a security test, an accidental overload, a temporary link problem,
and a routing fault are all controlled deviations from an expected baseline
that can be introduced, measured, and compared.

A first extension could support several families of adverse network conditions.
Examples include connection-pressure scenarios, controlled high-rate UDP or
ICMP traffic, service discovery and scanning inside the emulated network,
short traffic bursts, artificial congestion, temporary link failures, routing
misconfiguration, and selected routing anomalies. These examples represent
different causes, but the orchestration problem is similar in each case. The
framework must identify the participating nodes, validate that they belong to
the laboratory, apply a bounded action for a known period, continue measuring
the relevant services, and stop the action deterministically. The specific
traffic generator or fault mechanism would become an implementation detail of
a scenario type rather than the organising principle of the whole experiment.

This suggests separating three concepts that are only partially separated in
the current project: the \emph{topology}, the \emph{experiment}, and the
\emph{measurements}. The topology should continue to describe the relatively
stable network structure: nodes, interfaces, collision domains, addressing,
routing, and node resources. An experiment should describe a temporary
condition applied to that topology. A measurement specification should state
which observable properties must be recorded. Keeping these concepts separate
would allow the same topology to be reused for different experiments and the
same experiment definition to be applied, when compatible, to more than one
topology.

A possible future configuration could therefore use a dedicated experiment
document such as the conceptual example in
\Cref{lst:future-experiment-definition}. The syntax is illustrative and is not
implemented by the current framework.

\begin{lstlisting}[caption={Conceptual declarative definition of a controlled experiment.},label={lst:future-experiment-definition}]
experiment:
  name: service_resilience_test
  duration: 60

  source_nodes:
    - node_a
    - node_b

  target:
    node: server
    port: 9000

  workload:
    type: connection_stress
    profile: medium

  measurements:
    - availability
    - tcp_connect_time
    - application_response_time
\end{lstlisting}

The important aspect of this proposal is not the exact YAML syntax but the
change in abstraction. Instead of encoding an experiment primarily through
shell variables and script-specific command-line arguments, the user would
provide a structured description that could be validated before execution.
The framework could check that \texttt{node\_a}, \texttt{node\_b}, and
\texttt{server} are nodes in the selected topology, that the target port is
valid, that the requested workload type is installed, and that the selected
metrics are supported. The resulting experiment definition would then become
part of the reproducible input together with the topology configuration.

Scenario profiles could further reduce repetition. Rather than expressing every
parameter explicitly for every run, the framework could expose named profiles
such as \texttt{baseline}, \texttt{low}, \texttt{medium}, and \texttt{high}.
A profile would not imply an absolute notion of severity; it would simply be a
versioned collection of bounded parameters for a particular scenario type.
Users could then execute the same topology under multiple profiles while the
framework records the exact values used by each run. The current baseline and
controlled-load comparison in \Cref{chap:experimental-evaluation} can be seen
as a first manual instance of this more general pattern.

A generic experiment subsystem should also make safety constraints part of the
configuration model rather than relying only on conventions in individual
scripts. At minimum, future workload definitions should have explicit maximum
durations, maximum concurrency or rate parameters, and target restrictions.
Most importantly, a target should be expressed through a node reference in the
selected laboratory whenever possible, rather than as an arbitrary external
address. The validation phase could resolve the reference to the addresses
assigned to that node and reject destinations that do not belong to the
emulated environment. The framework could also terminate all experiment
processes when a configured deadline is reached or when the laboratory is
stopped. These constraints would preserve the intended use of the system as an
isolated experimentation platform.

Generalising the scenario model would also make mitigation testing more
systematic. The rate-limiting reverse-proxy experiment evaluated in
Section~\ref{sec:mitigation-results} is a first instance: a public endpoint at
\texttt{10.30.1.10} forwards admitted sessions to a private backend at
\texttt{10.30.2.10} while applying global and per-source connection limits.
The archived comparison reuses the same clients, round count, activation
schedule, and measurement schema across five executions per topology. These
paired repetitions were orchestrated as separate runs; future work should
encode their common parameters in a shared experiment definition and execute
the complete batch automatically.
More generally, a mitigation should not be evaluated by changing the workload
and the topology at the same time without recording those changes. The
framework could define experimental variants that share the same workload and
measurement specification while modifying one clearly identified network or
service configuration. Rate limiting, filtering, traffic shaping,
service-capacity changes, firewall policies, and architectural isolation are
examples of mechanisms that could be compared in this way.

Conceptually, the execution model would evolve from a scenario-specific script
into the pipeline shown in \Cref{fig:future-experiment-pipeline}.

\begin{figure}[htbp]
\centering
\begin{tikzpicture}[
    node distance=7mm,
    block/.style={draw, rounded corners, align=center, minimum width=4.8cm, minimum height=8mm},
    arrow/.style={-{Latex[length=2mm]}, thick}
]
\node[block] (topology) {Validated topology};
\node[block, below=of topology] (experiment) {Validated experiment definition};
\node[block, below=of experiment] (runtime) {Laboratory startup and readiness checks};
\node[block, below=of runtime] (scenario) {Bounded scenario execution};
\node[block, below=of scenario] (measure) {Concurrent measurements};
\node[block, below=of measure] (compare) {Scenario comparison and report};

\draw[arrow] (topology) -- (experiment);
\draw[arrow] (experiment) -- (runtime);
\draw[arrow] (runtime) -- (scenario);
\draw[arrow] (scenario) -- (measure);
\draw[arrow] (measure) -- (compare);
\end{tikzpicture}
\caption{Possible future pipeline for reusable controlled network experiments.}
\label{fig:future-experiment-pipeline}
\end{figure}

This extension would build directly on the present framework. The node names,
addressing information, output directories, concurrent client execution, and
CSV-based measurements already provide several of the required building
blocks. Future work would mainly need to formalise their relationships and
replace scenario-specific orchestration with a reusable experiment model.

\section{Runtime Verification and Automated Testing}
\label{sec:future-runtime-verification}

The current validation subsystem operates before laboratory generation. It can
identify malformed structures, duplicated addresses, inconsistent subnets,
invalid gateways, routing configuration problems, and other errors that can be
deduced from the declarative input. This is useful because it prevents many
problems from reaching the deployment stage. Nevertheless, a statically valid
configuration does not guarantee that the running network has reached the
state intended by the user. The difference between \emph{configuration
correctness} and \emph{runtime correctness} is therefore an important area for
future development.

Several causes can produce this difference. A container may fail to start, a
startup command may return an error, an expected interface may not be present,
a service process may terminate immediately, or a dynamic routing protocol may
require additional time to converge. Similarly, two routers can contain
syntactically valid BGP configurations while their session never reaches the
\texttt{Established} state. These conditions cannot be fully detected by
analysing the YAML document alone. A future framework should consequently add
a verification phase after \texttt{kathara lstart} and before an experiment is
considered ready.

The first level of runtime verification could compare the generated model with
the basic state of each container. The framework already knows which nodes
should exist and which interfaces and addresses were generated for them. It
could therefore verify that all expected containers are running, inspect
\texttt{ip addr} to check interface addressing, inspect \texttt{ip route} to
verify default and static routes, and confirm that required processes are
active. These checks would transform information already available during
generation into explicit post-start assertions.

Routing protocols require more specialised checks. For an OSPF laboratory, the
framework could inspect the neighbour table and verify that expected
adjacencies are present after a configurable convergence interval
\cite{rfc2328}. For BGP, it could check that selected peers reach the expected
session state and that required prefixes appear in the routing table
\cite{rfc4271,frr_user_guide}. The same principle applies to
application services: a node configured to provide a TCP service could be
considered ready only after the corresponding socket is listening and, where
appropriate, after a minimal application-level request succeeds. Such checks
would make startup deterministic from the point of view of later experiments.

The intended state could be partially inferred from the topology, but not every
expectation can be derived automatically. A future extension could therefore
allow users to declare selected behavioural expectations. The conceptual
configuration in \Cref{lst:future-expected-state} illustrates this idea.

\begin{lstlisting}[caption={Conceptual declaration of expected runtime properties.},label={lst:future-expected-state}]
expected:
  connectivity:
    - source: pc_a
      destination: server
      protocol: tcp
      port: 9000
      reachable: true

  routing:
    r1:
      ospf_neighbors: 2

    edge_router:
      bgp_peers_established: 1
\end{lstlisting}

The framework could translate such expectations into runtime checks and produce
a structured comparison between the expected and observed state. A failed
expectation should not necessarily be treated as a software error: in some
experiments, an unreachable destination may be intentional. The relevant point
is that the expected outcome is explicit and machine-readable, allowing the
same assertion to be evaluated across repeated runs.

This mechanism would also improve the orchestration logic. In the current
experimental workflow, fixed waiting periods and explicit reachability checks
are used to reduce the risk of starting measurements before the laboratory is
ready. A more general implementation could replace most fixed delays with
state-based readiness conditions. For example, instead of waiting an arbitrary
number of seconds after startup, the framework could poll the required OSPF
adjacencies until they appear or until a maximum convergence timeout is
reached. The experiment would start only when all mandatory readiness checks
succeed. If they do not, the run would be marked as invalid rather than being
mixed with application measurements. The distinction made in
\Cref{chap:experimental-evaluation} between a valid baseline and a run affected
by initial reachability problems shows why this separation matters.

Runtime verification should be complemented by automated testing of the
framework itself. The present implementation could be covered by unit tests for
configuration loading, structural validation, semantic validation, collision-
domain analysis, startup generation, \texttt{lab.conf} generation, and
FRRouting output. Each test could provide a small input configuration and
compare the resulting data or files with an expected representation. Negative
tests would verify that invalid configurations produce the correct validation
issues without modifying the output directory.

A second layer could consist of integration tests that create very small
Kathar\'a laboratories. An integration test could generate a two- or three-node
configuration, start it, execute one or more runtime assertions, and clean it
again. Such tests would verify not only the Python generation logic but also
its compatibility with the Kathar\'a and Docker environment. Because these
operations are more expensive than unit tests, they could be executed
separately or only on systems where the required runtime is available.

The combination of static validation, runtime verification, and automated
software tests would produce a clearer lifecycle:

\begin{center}
\texttt{configuration} $\rightarrow$ \texttt{static validation}
$\rightarrow$ \texttt{generation} $\rightarrow$ \texttt{startup}
$\rightarrow$ \texttt{runtime verification} $\rightarrow$
\texttt{experiment}
\end{center}

This evolution would also make the framework more useful in teaching. A
student laboratory could include both the intended topology and a set of
expected properties. The framework could then report whether the generated or
modified network satisfies those properties without hiding the underlying
Kathar\'a commands and files. The user would still be able to connect to nodes
and inspect the network manually, while automated checks would provide a
repeatable reference point.

\section{Experiment Management and Usability}
\label{sec:future-experiment-management}

The current project deliberately keeps its tools small and independent.
Generation, startup, connectivity measurements, metric visualisation, and
cleanup are implemented through separate Python or shell entry points. This
modularity was useful during development because each component could be tested
and changed independently. As the number of supported workflows grows,
however, the same modularity can create a usability problem: the user must know
which script to invoke, in which order, with which laboratory path, and where
the resulting files will be stored. A future version should preserve the
internal modules while presenting a more unified external interface.

A single command-line application could act as an orchestration layer. Its
interface could expose commands such as those shown in
\Cref{lst:future-cli}. The names are illustrative, but they express a consistent
resource-oriented workflow.

\begin{lstlisting}[language=bash,caption={Conceptual unified command-line interface.},label={lst:future-cli}]
framework validate topology.yml
framework generate topology.yml
framework start my_lab
framework verify my_lab
framework experiment my_lab experiment.yml
framework report my_lab experiment_0042
framework stop my_lab
framework delete my_lab
\end{lstlisting}

Internally, these commands could call the same generation, validation, testing,
and reporting components that exist as separate scripts. The change would
therefore concern orchestration rather than a monolithic rewrite. A unified
interface would also make help messages, error handling, configuration paths,
and naming conventions consistent across the project.

Experiment identity is another area that deserves stronger support. The current
repository already uses numbered experiment-result directories, which is an
effective first step because it prevents one run from overwriting the results
of another.
A future implementation could formalise this idea by assigning every execution
a unique experiment identifier and storing all inputs and metadata required to
interpret that run. The directory structure could resemble
\Cref{lst:future-results-layout}.

\begin{lstlisting}[caption={Possible future experiment output layout.},label={lst:future-results-layout}]
results/
`-- experiment_0042/
    |-- metadata.json
    |-- topology.yml
    |-- experiment.yml
    |-- measurements.csv
    |-- server_metrics.csv
    `-- report.pdf
\end{lstlisting}

The metadata file could record the experiment identifier, timestamp, framework
version, topology identifier, topology hash, experiment parameters, Kathar\'a
version, Docker version, and selected host information. Capturing this
information would make it easier to determine whether two result sets were
produced under comparable conditions. Storing a copy of the topology and
experiment definition alongside the measurements would also ensure that later
changes to files in \path{configs/} do not make an old result ambiguous.

A cryptographic hash of the input configuration could provide a simple link
between source and output. The hash would not replace the copied YAML document,
but it would make accidental configuration changes immediately visible. The
same principle could be extended to the framework version, for example by
recording a Git commit identifier when available. Reproducibility would then be
based not only on the numerical measurements but also on a traceable record of
how the laboratory and experiment were produced.

Reporting could also be integrated into the experiment lifecycle. The current
tools already calculate summary statistics and visualise selected measurements.
A future report generator could systematically combine raw measurements,
metadata, validation status, runtime readiness checks, and plots. It could
produce an HTML or PDF report containing success rate, availability, mean,
median, percentile statistics, error categories, resource utilisation, and
comparisons among profiles or experiment variants. The raw CSV files should
remain available so that the automated report is an additional representation
rather than the only form of the data.

A related improvement is automatic topology visualisation. The generator
already has enough information to construct a graph: node names, node types,
interfaces, collision domains, addresses, and routing relationships are known
before the laboratory is started. A visualisation component could transform
this model into Graphviz, Mermaid, SVG, or another portable representation. The
same diagram could be embedded in generated experiment reports and used as a
quick check before deployment. For larger laboratories, automatic layout would
be especially useful because a textual YAML file becomes progressively harder
to interpret as the number of nodes and collision domains increases.

Visualisation should remain derived from the declarative source rather than
becoming a second authoritative configuration. An editor that allows users to
modify the graph directly could be considered later, but the underlying change
should still be represented in the same data model and serialised back into a
versioned configuration. This avoids creating two independent descriptions of
the same laboratory.

A lightweight web interface could eventually expose generation, validation,
lifecycle control, topology visualisation, and result browsing through a
browser. Such an interface could be useful in teaching environments where
users are less familiar with command-line tools. Nevertheless, the CLI should
remain a first-class interface because it is more suitable for automation,
version-controlled workflows, remote execution, and integration testing. A web
interface would therefore be a front end to the same core services rather than
a replacement for them.

Finally, experiment management could support batches of related runs. Instead
of manually executing one profile at a time, the user could declare a matrix
of scenarios, repetitions, or topology variants. The framework would assign a
separate identifier to each run, execute them sequentially or with explicitly
controlled parallelism, and generate an aggregate report. This would be
particularly useful when studying variability. The repeated-run method used in
\Cref{chap:experimental-evaluation} already bases its principal conclusions on
a distribution of run-level summaries, but its five acquisition blocks were
launched as separate commands. A batch facility would make the same method
declarative, automatic, and reusable for other experiment matrices.


\section{Using LLM-Based Coding Alongside the Framework}
\label{sec:future-llm-integration}

The comparison in \Cref{sec:llm-role} first suggests an immediately usable
workflow rather than a requirement to integrate a model into the software. A
user can employ any general-purpose LLM-based coding assistant to draft or
modify the YAML file, run the existing validator, and return its diagnostics to
the assistant when another revision is useful. The model does not write
directly to the Kathar\'a runtime and is not treated as the authority on whether
a configuration is valid: its output enters the existing pipeline as untrusted
candidate input. A future front end could automate this exchange, but the main
synergy already exists through ordinary files, command output, and measured
results.

The first implementation step would be to formalise the currently implemented
configuration model. The current accepted language is defined by Python code
and documented in \Cref{chap:configuration-validation}; future work already identifies a strict
schema as desirable. An LLM integration strengthens that requirement because a
machine-generated draft benefits from an explicit description of allowed keys,
types, enumerations, and nested structures. A JSON Schema or an equivalent
versioned schema could be supplied to the model as context and also enforced by
the framework independently. The model could then be instructed to produce
only a structured object conforming to that schema instead of free-form
Kathar\'a files.

Schema conformance would remain only the first check. The existing semantic
validator would still be needed because many network properties depend on
relationships rather than local field types. Two addresses can each be valid
CIDR values while still being duplicates; a gateway can be a valid IP address
while being unreachable from the node; and two BGP declarations can each be
well formed while disagreeing on the remote autonomous system. These are
precisely the kinds of checks for which deterministic code is preferable to
asking the same model that produced the configuration whether its own answer is
correct.

The resulting interaction could use the validation report as a repair protocol.
A first LLM response would produce a candidate YAML document. The framework
would return either success or a list of structured issues. If errors were
present, a controller could send the candidate configuration and the relevant
issue codes back to the model and request a revised version. The loop would be
bounded by a small maximum number of attempts and would preserve every candidate
for traceability. A conceptual exchange is shown in
\Cref{lst:future-llm-loop}.

\begin{lstlisting}[language=bash,caption={Conceptual LLM-assisted validation loop.},label={lst:future-llm-loop}]
user intent
  -> LLM: candidate topology.yml
  -> framework validate topology.yml
       -> VALID: generate laboratory
       -> ERROR: return structured diagnostics
                 -> LLM/human revises topology.yml
                 -> validate again
\end{lstlisting}

This design would make the LLM replace part of the manual authoring effort, not
the validator. It would also preserve a useful property of the present project:
a user who does not want to use an LLM could continue to write YAML directly
and follow exactly the same validation and generation path. The AI component
would therefore be an adapter to the declarative interface rather than a new
required dependency.

A second use of LLMs could occur after execution. Runtime verification proposed
in \Cref{sec:future-runtime-verification} could produce structured facts such as
missing routes, inactive services, failed OSPF adjacencies, or unexpected
application responses. An LLM could combine those facts with the source
configuration and selected logs to explain likely causes or suggest a next
inspection command. The same principle applies to experiment results: raw CSV
files, packet captures, and metadata should remain authoritative evidence, while
an LLM may generate a human-readable interpretation or propose hypotheses for a
subsequent experiment. This separation avoids making an automatically generated
summary the only record of what happened.

The integration also introduces security and control requirements. The present
configuration language intentionally exposes powerful escape mechanisms. A
\texttt{commands} entry becomes shell text in a startup file,
\texttt{frr.config} can contain arbitrary FRRouting configuration, and an
explicit image can change the software executed by a node. These fields are
useful for expert users but should not be executed automatically merely because
an LLM proposed them. A safer assisted mode could prohibit opaque fields by
default, require explicit human approval for them, or apply allow-lists for
commands and images. Generation could also run under resource limits and with a
clearly isolated output directory before any laboratory is started.

Traceability should include the AI-assisted part of the process. If an LLM is
used, experiment metadata could record that fact together with the model family
or endpoint, the prompt template or task description, the schema version, and a
hash of the accepted YAML source. The exact conversational history need not
become part of every experiment, but the final source and framework version
must remain sufficient to regenerate the laboratory. This preserves the thesis'
main reproducibility principle: the accepted configuration, not the transient
conversation, is the authoritative experiment input.

There is also a useful teaching perspective. An LLM can explain why a route or
BGP relationship is required and can help a student produce an initial
configuration, while deterministic validation prevents the interaction from
ending at a plausible-looking answer. The student can compare the LLM proposal,
the validator diagnostics, the generated Kathar\'a files, and the observed
runtime state. The AI assistant then becomes part of the learning process
without hiding the concrete networking mechanisms that the laboratory is meant
to expose.

Research on LLM-assisted networking already points toward such hybrid designs.
NetConfEval considers the interaction between LLMs and existing network
synthesis mechanisms, while LLM-NetCFG combines intent-driven generation with
verification in a network-management workflow
\cite{wang2024netconfeval,lira2025llmnetcfg}. The proposed extension follows
the same general principle but keeps the scope specific to this thesis: the LLM
would improve how a user reaches the YAML representation and how diagnostics are
interpreted, whereas the framework would continue to own validation, generation,
and the transition to the Kathar\'a runtime.

\section[Towards a More General Network Emulation Framework]{Towards a More General Network Emulation\\Framework}
\label{sec:future-general-framework}

The previous proposals can be implemented incrementally without abandoning
Kathar\'a. In fact, preserving direct compatibility with Kathar\'a is one of the
main strengths of the current project: generated laboratories remain standard,
inspectable environments rather than opaque artefacts that require a custom
runtime. If the framework grows substantially, however, a stronger separation
between its logical network model and its Kathar\'a-specific generation code
would improve maintainability and make the architecture easier to extend.

The current pipeline can be simplified as

\begin{center}
\texttt{YAML configuration} $\rightarrow$ \texttt{Kathar\'a files}.
\end{center}

A more general architecture would insert an explicit internal model and a
backend boundary. The existence of different emulation approaches, including
Kathar\'a and container-/namespace-based systems such as Mininet, illustrates
why the logical model and backend-specific representation should remain
separate \cite{scazzariello2020kathara,lantz2010network}:

\begin{center}
\texttt{configuration} $\rightarrow$ \texttt{network model}
$\rightarrow$ \texttt{backend} $\rightarrow$ \texttt{runtime artefacts}.
\end{center}

The internal model would represent concepts such as laboratories, nodes,
interfaces, collision domains, addresses, routes, and routing configuration
without embedding the syntax of \texttt{lab.conf} or startup files. The
Kathar\'a backend would then translate this model into the artefacts described
in \Cref{chap:laboratory-generation}. This separation would make it easier to
test the model independently from file generation and would provide a clear
place for future backend-specific behaviour.

Such an abstraction does not imply that support for other emulation platforms
would be simple. Different systems expose different network models, lifecycle
semantics, privilege requirements, and configuration mechanisms. A future
backend interface would therefore be valuable primarily because it separates
responsibilities, not because it guarantees portability. Kathar\'a would remain
the reference backend, while alternative targets could be investigated only
where their capabilities map cleanly onto the internal model.

A plugin architecture could complement this separation. The present
\texttt{gen\_lab.py} module contains several responsibilities related to
routing, resource import, Wireshark integration, and output generation. If more
protocols, experiment types, validators, and reporting components are added,
continuing to extend a small number of central modules would increase coupling.
A future layout could instead organise optional capabilities around explicit
extension points, for example:

\begin{lstlisting}[caption={Conceptual modular organisation for a larger framework.},label={lst:future-plugin-layout}]
framework/
|-- core/
|-- model/
|-- validators/
|-- backends/
|   `-- kathara/
|-- experiments/
|-- measurements/
|-- reporting/
`-- plugins/
    |-- ospf/
    |-- bgp/
    `-- wireshark/
\end{lstlisting}

A plugin would not necessarily correspond to a dynamically installed Python
package. The important design principle is that an optional capability should
have a defined way to contribute configuration fields, validation rules,
generation logic, runtime checks, and documentation. For example, a routing
extension could register both semantic validation and backend generation rules
without requiring unrelated experiment-management code to know its internal
details.

Chapter~\ref{chap:robustness-security-scalability} already characterises the
scalability of the translator with synthetic inputs, including configurations
containing up to 10,000 nodes. Those measurements cover loading, semantic
validation, and file generation, but deliberately stop before the generated
laboratories are started. The remaining scalability question is therefore
runtime execution: future work should measure Kathar\'a startup time, host
memory and CPU consumption, container stability, and routing convergence as the
number of executed nodes and collision domains increases. This would separate
the scalability of the Python translation stage, which has been benchmarked,
from the scalability of the Docker/Kathar\'a runtime, which has not.

Large experiments also raise questions about orchestration. A single host is a
simple and reproducible environment for the present project, but CPU scheduling,
memory pressure, and the number of containers eventually influence results.
The framework should therefore record host resource information and make clear
when measurements are specific to one execution environment. Distributed
execution could be studied in the future, but it would represent a substantial
change in scope because synchronisation, clock behaviour, network links between
hosts, and result aggregation would become part of the experimental model.
For this reason, distributed operation should be considered a possible
long-term direction rather than a direct next step.

The same caution applies to a graphical topology editor and other usability
features. They can make the framework easier to approach, but they should not
weaken the properties that make the current design useful for reproducible
experiments. The declarative source must remain versionable, inspectable, and
sufficient to reconstruct the laboratory. Runtime state and generated files
should remain derived artefacts. Any additional interface should therefore
operate on the same underlying model rather than creating an independent source
of configuration.

Taken together, the developments discussed in this chapter describe a gradual
path from a laboratory generator towards a more complete experimentation
framework. The first stage would generalise scenario execution and runtime
verification using the existing Kathar\'a backend. The second would improve
experiment identity, metadata, reporting, and automated testing. An optional
LLM-assisted front end could then use the same declarative model and validation
report without becoming a replacement for them. Only after these foundations
are stable would a richer internal model, plugin system, or additional backends
be necessary. This staged evolution preserves the simplicity of the present
implementation while providing clear architectural directions for extending its
use in network emulation, resilience testing, security experimentation,
AI-assisted configuration, and teaching.
